\section{Exponential maps}
Let $M$ be a Riemannian manifold for $p \in M$ and $v \in T_p M$, let $\gamma_v : J_v \to M$ be the maximal geodesic satisfying
$\gamma_v(0)  =p$ and $\gamma_v'(0) = v$, where $J_v$ is an open interval containing $0$.

\begin{definition}[Exponential map]
For $v \in TM$ if $1 \in J_v$, define $\exp v := \gamma_v(1)$.
Let $D:= {v \in TM | 1 \in J_v}$.
The map 
\[
\exp : D \to M
\]
is the \emph{exponential map}.
For $p \in M$, let $D_p := D \cap T_pM$. The map
\[
\begin{aligned}
    &\exp_p : D_p \to M \\
    &\exp_p := \exp|_{D_p},
\end{aligned}
\]
is the \emph{exponential map} at $p$.

$\exp_p v \quad (v \in T_pM)$ is the point reached by traveling along the geodesic from $p$ in direction $v$ for the distance $\|v\|$.
$\forall v \in D_p$ and $0 \leq \forall t \leq 1$ $1 \in J_v = t J_{tv} \subset J_{tv} \implies tv \in D_p$.
Thus, $D_p$ is a star shaped domain.
\end{definition}

\begin{proposition}
    $D \subset TM$ and $D_p \subset T_p M$ are open.
    Both $\exp$ and $\exp_p$ are $C^{\infty}$
    \label{prop:10.2}
\end{proposition}
\begin{proof}
Follows from \autoref{prop:7.10}.
\end{proof}

Let $0_p$ denote the zero vector in $T_p M$.
\begin{proposition}[The differential map]
    $d(\exp_p)_{0_p} : T_{0_p} T_p M \approx T_p M \to T_p M$ is the identity.
    \label{prop:10.3}
\end{proposition}

\begin{proof}
    For $v \in T_p M$, let $c(t) := tv$. Then $(\exp_p \circ c)(t) = \exp_p(tv) = \gamma_v(t)$.
    This shows that: 
    \[
    \begin{aligned}
        d \exp_p(v) &= d \exp_p (c'(0))= (\exp_p \circ c)'(0) \\
        &= \gamma_v'(0) = v 
    \end{aligned}
    \]
\end{proof}

Let $\pi : TM \to M$ be the natural projection.

\begin{lemma}
    Define $f: D \to M \times M$ by
    \[
    \begin{aligned}
        \exp f(v) := \left( \pi(v), \exp v\right), v \in D
    \end{aligned}
    \]
    $\forall p_0 \in M$, $\exists$ open neighbourhood $0_{p_0} \subset TM$ of $0_{p_0}$ such that
    \begin{enumerate}
        \item $0_{p_0} \subset D$
        \item $\exp f(0_{p_0})$ is open in $M \times M$
        \item $\exp f : 0_{p_0} \to \exp f(0_{p_0})$ is a diffeo.
    \end{enumerate}
    \label{lemma:10.4}
\end{lemma}

\begin{proof}
    Fix chart $(U, \phi)$ aroun $p_0$ $\implies$ induced chart $(TU,\tilde{\phi})$ for $TM$. let:
    $\tilde{\phi}(v) = (x^1, \dots, x^n, v^1, \dots, v^n)$.
    $\phi(\exp v) = (y^1, \dots, y^n)$.
    Local expression $\bar{f} := (\phi \times \phi) \circ \exp f \circ \tilde{\phi}^{-1}$ is:
    \[
    (f^1, \dots, f^{2n}) := \bar{f}(x^1, \dots, x^n, v^1, \dots, v^n) = (x^1, \dots, x^n, y^1, \dots, y^n)
    \]
    Jacobian of $\bar{f}$ at $(\phi(p_0),0)$ is non-singular.
    By Inverse Function Theorem, $\exists$ open $0 \subset \mathbb{R}^{2n}$ such that
    $\bar{f}|_0$ is a diffeo and $\bar{f}(0)$ is open/
    Set $0_{p_0} := \tilde{\phi}^{-1}(0)$
\end{proof}

\begin{problem}
    Provide details for Proof of \autoref{lemma:10.4}.
    Specifically, show non-singularity of Jacobian of $\bar{f}$ at 
    $(x^1, \dots, x^nm v^1, \dots, v^n) = (\phi(p_0),0. \dots,0)$. 
    \label{problem:10.5}
\end{problem}

\begin{lemma}\label{lemma:10.5}
    $\forall p_0 \in M$, $\exists$ open neighbourhood $0$ of $p_0, \quad \exists r >0$ such that $\forall p \in 0$:
    \begin{enumerate}
        \item $B(0_p,r) \subset D_p$
        \item $\exp_p (B(0_p,r))$ is open in $M$
        \item $\exp_p : B(0_p,r) \to \exp_p(B(0_p,r))$ is a diffeo.
    \end{enumerate}
    Here $B(0_p,r) \subset T_p M$ is the open ball of radius $r$ centered at $0_p$.
\end{lemma}
\begin{proof}
    Let $f: 0_{p_0} \to f(0_{p_0})$ be the diffeo from \autoref{lemma:10.4}, since $0_{p_0}$ is an open neighbourhood of $0_{p_0}$,
    $\exists O \subset M$: open around $p_0$ such that $r>0$ such that
    \[
    TO(r) := {v \in TM | \pi(v) \in O, |v| <r} \subset 0_{p_0}
    \]
    $\forall p \in O$, $B(0_p,r) \subset TO(r) \subset 0_{p_0} \subset D \implies (1)$ holds.
    Define $Q_P := {q \in M| (p,q) \in f(0_{p_0})}$
    Since $f(0_{p_0}) \subset M \times M$ is open, $Q_p$ is also open.

    \[
\left(
\begin{array}{c@{\;}c@{\;}c@{\;}c}
    h \colon & M & \to & M \times M \\
    & \rotatebox[origin=c]{90}{$\in$} & & \rotatebox[origin=c]{90}{$\in$} \\
    & q & \mapsto & (p, q)
\end{array}
\text{ is continuous and } Q_p = h^{-1}(f(O_{p_0})).
\right)
\]
The inverse
\[
f^{-1}(p,q)  = \exp_p^{-1}(q), \quad (p,q) \in f(0_{p_0})
\]
is $C^{\infty}$, so $\exp_p^{-1}\|_{Q_p} : Q_p \to T_p M \cap 0_{p_0}$ is $C^{\infty}$ for fixed $p$.
Thus $\exp_p: T_p M \cap 0_{p_0} \to Q_p$ is a diffeomorphism. Since $B(0_p,r) \subset 0_{p_0}$, (2) and (3) follow.
\end{proof}

\begin{theorem}[Gauss Lemma]
    $\forall p \in M$ and $\forall v,w \in T_pM$, the following hold:
    \begin{enumerate}
        \item $d(\exp_p)_v v = \gamma_v'(1)$
        \item $\langle \gamma_v'(1), d(\exp_p)_v w \rangle = \langle v, w \rangle$.
    \end{enumerate}
    Here, $w$ is identified with a vector at $v \in T_pM$ via $T_pM \approx T_v T_pM$
    \label{theorem:10.6}
\end{theorem}

\begin{proof}
\begin{enumerate}[label=(\arabic*)]
\item follows from $\gamma_v'(1) = \frac{d}{dt} \exp_p (tv) |_{t=1}= d(\exp_p)_v v$.

\item Let $ p \in M$ and $v,w \in T_pM$. Case $v=0_p$ is trivial.
Assume $ v \neq 0_p$. consider variation:
\[
\begin{aligned}
&c_s (t) := c(s,t) := \exp_p (t(v+sw)), \quad s \in (- \epsilon, \epsilon), 0 \leq t \leq 1. \\
&c_0(t) = \exp_p (tv) = \gamma_v (t) \text{ is a geodesic}
\end{aligned}
\]
$\implies \nabla_{c_0} = 0$. Variation field $X(t) = d(\exp_p)$.
By 1st variation formula (cf. \autoref{theorem:7.2})
\[
\begin{aligned}
\frac{d}{ds} L(C_s)|_{s=0} &= \frac{1}{|c_0'(0)|} \left( \langle c_0'(1, X(1)) \rangle - \langle c_0'(0), X(0) (=0) \rangle \right) \\
&= \frac{1}{|v|} \langle \gamma_v'(1), d(\exp_p)_v(w) \rangle
\end{aligned}
\]
Alternativly, $L(Cs) = |v+sw| = \langle v+sw, v +sw \rangle^{1/2}$
\[
\frac{d}{ds} L(Cs) |_{s=0} = \frac{\langle \frac{d}{ds}(v+sw), v +sw \rangle|_{s=0}}{|v|} = \frac{\langle v,w \rangle}{|v|}
\]
\end{enumerate}
\end{proof}

\begin{remark}
    Combining (1) and (2) yields:
    $$
    \langle d(\exp_p)_v v, d(\exp_p)_v w \rangle = \langle v,w \rangle
    $$
    Thus, Gauss Lemma is weaker than $d(\exp_p)v$ being isometric.
    If ${v,w}$ is linearly dependent $|d(\exp_p)_v w | = |w|$.
    In general $d(\exp_p)_v$ is not isometric.
    For example, $M = \mathbb{S}^2 \subset \mathbb{R}^3, \quad v \in T_pM; \|v\|=\pi$.
    $d(\exp_p)_v$ is an isometry $\implies$ it is injective.
    By the inverse function theorem, $\exp_p$ is a diffeo on a neighbourhood of $v$. However, this is impossible since for $\forall v' \in T_pM$ satisfying 
    $\|v'\| = \pi, \quad \exp_p v'$ coincides with the antipodal point of $p$.
    \label{remark:10.7}
\end{remark}

\begin{lemma}\label{lemma:10.8}
$\forall p \in M$, $\forall v \in T_pM$, let $\sigma : [0,l]\to D_p$ be a piecewise smooth regular curve s.t. $\sigma(0)  0_p$ and
$\sigma(l)= v$. Then
\[
L(\exp_p \odot sigma) \geq ||v||.
\]
furthemore, if $d(\exp_p)_{\sigma(t)}$ is non-singular $\forall t \in [0,l]$, then
for every monotonically increasing function $f:[0,l] \to [0,1]$ with $f(0) = 0, f(l) = 1$,
\[
\sigma(t) = f(t)v.
\]
\end{lemma}

\begin{proof}
Assume $v \neq 0$. Let 
\[
t_0 := \sup\{t \in [0,l] : \sigma(t) = 0_p\}.
\]
$\forall t \in (t_0, l)$, let $r(t) := ||\sigma(t)|| >0$ and $e(t) \in t_pM$ be the unit vector s.t. $\sigma(t) = r(t)e(t)$.
Let also $c(t):= \exp_p \odot \sigma(t)$. Since $\sigma'(t) = r'(t)e(t) + r(t)e'(t)$, 
\[
\begin{aligned}
c'(t) &= d(\exp_p)_{\sigma(t)}\sigma'(t) \\
&= r'(t) d(\exp_p)_{\sigma(t)}e(t) + r(t)d(\exp_p)_{\sigma(t)}e'(t) \\
&= \frac{r'(t)}{r(t)}\gamma'_{\sigma(t)}(1) + r(t)d(\exp_p)_{\sigma(t)}e'(t).
\end{aligned}
\]
By Gauss' lemma,
\[
\begin{aligned}
\langle \gamma'_{\sigma(t)}(1),d(\exp_p)_{\sigma(t)}e'(t)\rangle &= r(t) \langle e(t), e'(t)\rangle \\
&= \frac{r(t)}{2} \frac{d}{dt} \langle e(t), e(t) \rangle = 0.
\end{aligned}
\]
Thus the decomposition of $c'(t)$ above is an orthogonal decomposition.
Since $|\gamma'_{\sigma(t)}(1)| = ||\sigma(t)|| = r(t)$, we have
\[
|c'(t)| = \sqrt{|r'(t)|^2 + ||r(t)d(\exp_p)_{\sigma(t)}e'(t)||^2} \geq |r'(t)| \geq r'(t).
\]
Integrating gives
\[
\begin{aligned}
L(c|_{[t_0,l]}) &= \int_{t_0}^l|c'(t)|dt \geq \int_{t_0}^lr'(t)dt \\
&= ||\sigma(l)|| - ||\sigma(t_0)|| = ||\sigma(l)|| = ||v||.
\end{aligned}
\]
This proves $L(\exp_p\odot\sigma) \geq ||v||$.

Now assume $d(\exp_p\odot\sigma)$ is non-singular and $L(\exp_p\odot\sigma)= ||v||$. Then the above inequality implies that 
$L(\exp_p\odot\sigma|_{[0,t_0]}) = 0 \Rightarrow \sigma(t) = 0_p$ on $[0,t_0]$.
This implies
\[
|c'(t)| = |r'(t)| = r'(t).
\]
Thus, the second term in the othogonal expansion of $c'(t)$ must vanish, i.e. $d(\exp_p)_{\sigma(t)}e'(t) = 0$.
Since $d(\exp_p)_{\sigma(t)}$ is non-singular, $e'(t) = 0_p$, implying that $e(t)$ is constant. So
\[
e(t) = e(l) = v/||v||.
\]
Setting $f(t) := r(t)/||v||$ yields $\sigma(t)=f(t)v$.
\end{proof}

\begin{lemma}\label{lemma:10.9}
Let $p \in M$, $r>0$. Assume $B(0_p,r) \subset D_p$ and $\exp_p : B(0_p,r) \to \exp_p(B(0_p,r))$ is a diffeoorphism.
Then
\begin{enumerate}[label=(\arabic*)]
\item $\forall q \in \exp_p(B(0_p,r))$, $\exists ! v \in B(0_p,r)$ s.t. $q = \exp_pv$.
Any piecewise smooth minimizing curve connecting $p$ and $q$ coincides with the geodesic $\gamma_v : [0,1] \to <$ up to
reparametrization.
\item $\exp_p(B(0_p,r)) = B(p,r)$.
\end{enumerate}
\end{lemma}

\begin{proof}
\begin{enumerate}[label=(\arabic*)]
\item Existence and uniqueness of $v$ follows from the bijectivity of $\exp_p$.
Let $c:[0,1] \to M$ be a piecewise smooth curve from $p$ to $q$. Define
\[
\sigma(t) := (\exp_p|_{B(0_p,r)})^{-1}\odot c(t),
\]
whenever $c(t) \in \exp_p(B(0_p,r))$.
\begin{enumerate}[label=(\roman*)]
\item If $c([0,1]) \not\subset \exp_p(B(0_p,r))$: Let $0<\rho < r$.
Define $l := \sup\{t\in[0,1]:c([0,t)\subset\exp_p(B(0_p,\rho))\}$.
set $U := \exp_p(B(0_p,\rho))$ and $\phi := (\exp_p|_{B(0_p,\rho)})^{-1}:U \to B(0_p,\rho)$ and
identiy $T_pM \approx \mathbb{R}^n$. Applying \autoref{lemma:1.16} to $(U,\phi)$ it follows that
$\sigma(l) \in \partial B(0_p,\rho)$ since $\sigma(l) \notin B0_p,\rho)$. % I think Lemma 1.16 in his notes is 1.23 in our notes.
By \autoref{lemma:10.8}, we have 
\[
L(c) \geq L(c|_{[0,l]}) \geq ||\sigma(l)|| = \rho.
\]
Since $\rho$ was arbitrary, $L(c) \geq r$. However, $L(\gamma_v|_{[0,1]}) = ||v|| < r$. Thus,
\[
L(\gamma_v)_{[0,1]} < r \leq L(c).
\]
\item If $c([0,1])\subset\exp_p(B(0_p,r))$: $\sigma(t)$ is defined $\forall t \in [0,1]$.
By \autoref{lemma:10.8}, $L(c) \geq ||v|| = L(\gamma_v|_{[0,1]})$.
\end{enumerate}
In both cases, $L(c) \geq L(\gamma_v|_{[0,1]})$, so $\gamma_v$ is a minimizing curve
connecting $p$ and $q$. Then, $L(\gamma_v|_{[0,1]}) = L(c)$. In this case, if $c([0,1]) \not\subset\exp_p(B(0_p,r))$, then by
(i), we have $L(\gamma_v)_{[0,1]} < L(c)$, which is a contradiction. Therefore, $c([0,1]) \subset \exp_p(B(0_p,r))$. 
Furthermore, by the last part of \autoref{lemma:10.8}, there exists a non-decreasing function $f:[0,1]\to[0,1]$ s.t.
\[
f(0)=0,\quad f(1)=1, \quad \sigma(t) = f(t)v\quad \text{ for } t \in [0,1].
\]
Thus, by reparametrizing $c$ as $\tilde{c}(t) := (c\odot f^{-1})$, we have 
\[
\tilde{c}(t) = c(f^{-1}(t)) = \exp_p(f(f^{-1}(t))v) = \exp_p(tv) = \gamma_v(t).
\]
\item For $q \in \exp_p(B(0_p,r))$, let $q=\exp_p$.
Since $\gamma_v$ is minimizing,
\[
d(p,q) = L(\gamma_v) = ||v|| < r \Rightarrow q \in B(p,r).
\]
Conversely, let $q \in B(p,r) \Rightarrow d(p,q) < r$. Then, there exists a curve $c$ connecting $p$ and $q$ s.t. $L(c) <r$.
If $c([0,1]) \not\subset \exp_p(B(0_p,r))$, by (1i), $r \leq r$, a contradiction.
Thus $c([0,1]) \subset \exp_p(B(0_p,r))$, specifically $q \in \exp_p(B(0_p,r))$.
\end{enumerate}
\end{proof}

Combining \autoref{lemma:10.5} and \autoref{lemma:10.9}, we have the following:

\begin{theorem}\label{theorem:10.10}
$\forall p_0 \in M$, there exists an open neighbourhood $0 \subset M$ of $p_0$ and $r >0$ s.t. the following hold
\begin{enumerate}[label=(\arabic*)]
\item $\forall p \in 0$, $B(0_p,r) \subset D_p$ and $\exp_p:B(0_p,r)\to B(p,r)$ is a diffeomorphism.
\item $\forall p \in 0$ and $\forall q \in B(p,r)$, there exists a unique piecewise smooth minimizing curve
(up to renormalization) connecting $p$ and $q$. Thus curve is the geodesic $\gamma_v:[0,1]\to M$ where
\[
v = (\exp_p|_{B(0_p,r)})^{-1}(q) \in B(0_p,r).
\]
\end{enumerate}
\end{theorem}
A piecewise smooth regular curve $c : (a,b) \to M$ is locally minimizing if and only if $\forall t_0 \in (a,b)$,
$\exists \delta > 0$ s.t. $[t_0 - \delta, t_0 + \delta] \subset (a,b)$ and $c|_{[t_0 - \delta, t_0 + \delta]}$
is the minimizing curve between $c(t_0 - \delta)$ and $c(t_0 + \delta)$.
Sub-arcs of minimizing curves are themselves minimizing, in particular every minimizing curve is locally minimizing (verify this yourself).

\begin{theorem}\label{theorem:10.11}
Let $\gamma:(a,b) \to M$ be a piecewise smooth regular curve. Then
\[
\gamma \text{ is a geodesic} \Leftrightarrow \gamma \text{ is locally minimizing}.
\]
\end{theorem}

\begin{proof}
($\Rightarrow$): $\forall t_0 \in (a,b)$, let $p_0:=\gamma(t_0)$. Apply \autoref{theorem:10.10} to get $0$ and $r>0$.
Choose $\delta >0$ s.t. $\gamma < r/2$, $[t_0-\delta,t_0+\delta]\subset(a,b)$ ane $\gamma(t_0-\delta) \in 0$.
Since $L(\gamma|_{[t_0-\delta,t_0+\delta]})=2\delta < r$, $\gamma([t_0-\delta,t_0+\delta])\subset B(\gamma(t_0-\delta),r)$.
By \autoref{theorem:10.10}, for $p:=\gamma(t_0-\delta)$ and $q := \gamma(t_0+\delta)$, the curve $\gamma|_{[t_0-\delta,t_0+\delta]}$ is minimizing.
Thus $\gamma$ is locally minimizing.

($\Leftarrow$): Let $t_0 \in (a,b)$, $\delta' >0$ s.t. $\gamma|_{[t_0-\delta',t_0+2\delta']}$ is minimizing.
Apply \autoref{theorem:10.10} at $p_0 := \gamma(t_0)$ to get $0$ and $r>0$.
Choose $\delta >0$ s.t. $\delta \leq \delta'$, $\delta < r/2$ and $r(t_0 - \delta) \in 0$.
let $p:= \gamma(t_0 - \delta)$, $g := \gamma(t_0 + \delta)$. Since $\delta \leq \delta'$,
$\gamma|_{[t_0-\delta,t_0+\delta]}$ is minimizing between $p$ and $q$. since $d(p,q) = d\delta <r$, \autoref{theorem:10.10} (2) implies that
$\gamma|_{[t_0-\delta,t_0+\delta]}$ is a geodesic.
Since $t_0$ was arbitrary, $\gamma$ is a geodesic.
\end{proof}

\begin{definition}[Normal neighbourhood and coordinates]
assume $\exp_p:B(0_p,r) \to B(p,r)$ is a diffeomorphism for $p \in M$ and $r >0$. Let $\{e_i\}$ be an 
orthonormal basis of $T_pM$. For $q \in B(p,r)$, define
\[
x^i := \langle (\exp_p|_{B(0_p,r)})^{-1},e_i\rangle, \qquad (i=1,\dots,n),
\]
so that $q = \exp_p\left(\sum_{i=1}^n x^i e_i\right)$. The coordinates $\phi(q) = (x^1,\dots,x^n)$ defines a coordinate chart $(B(p,r),\phi)$.
this is the \emph{normal coordinate neighbourhood} at $p$ subordinated to $\{e_i\}$ and $(x^1,\dots,x^n)$ are \emph{normal coordinates}.
\end{definition}

The following is a consequence of \autoref{theorem:10.10}:

\begin{corollary}
$\forall p \in M$, normal coordinate neighbourhood exists around $p$.
\end{corollary}

Let $\{e_i\}$ be an orthonormal basis of $T_pM$, $(x^1,\dots,x^n)$, the subordinated normal coordinates at $p$ and $E_i := \partial/\partial x^i$.
\begin{problem}\label{problem:16}
Show that $\forall i,j,k = 1,\dots,n$:
\begin{enumerate}[label=(\arabic*)]
\item $(E_i)_p = e_i$,
\item $\nabla_{e_i}E_j = 0$ and $\Gamma_{ij}^k(p) = 0$.
\end{enumerate}
\end{problem}

\begin{remark}
The equlities is \autoref{problem:16} hold only at a specific point $p$.
\end{remark}